\section{Introduction}

World models often use a fixed observation window. An observation can affect
the future after it leaves this window, but the model can no longer attend to
it directly. Extending the window increases cost and does not decide which
past observations are useful. A practical external memory must instead answer
three questions: when to write, what to retain, and what to retrieve for a
particular prediction.

We study this problem without supplying event times or event labels.
Causal-Effect Memory (\cem) uses the prediction loss of a fixed world model as
its training and evaluation signal. Prediction surprise proposes a visual
event. A delayed test measures whether retaining the event reduces
action-conditioned future-latent loss. A learned score approximates this
measurement, and periodic group deletion updates the score. At inference, a
bounded router ranks verified events using query content, event content, age,
and predicted need.

Our main experiment uses original OGBench renderings. The model receives only
frames, actions, and timestamps. Frozen DINOv2 patch tokens provide visual
features, and an action-conditioned latent predictor is trained on each
environment and then frozen. These predictors are \emph{DINO-feature
action-conditioned world models}; they follow the feature-space design of
DINO-WM \citep{r_dinowm}, but they are not released DINO-WM checkpoints.
The study covers 10 environments from navigation, manipulation, puzzle, and
scene families.

The raw results are positive but limited. \cem\ lowers four-step prediction
error relative to no-memory, reset, shuffled-episode, and matched-random
controls in every environment, and the improvement grows with rollout
horizon. A recent-only extension is nevertheless slightly better on average.
Moreover, the learned causal-effect ranking has low held-out correlation with
measured deletion effects. Thus past state is useful, but the current method
does not reliably identify the most useful event.

Two controlled studies clarify this result. First, the released frozen
DINO-WM Wall host shows a large gap between selected-event and random-event
deletion under a separate encoded event-versioning protocol. Second, a frozen
LeWorldModel (LeWM) \citep{r_lewm} can retain a six-way binding in memory while
failing to use it without harming its original prediction loss. We call
\emph{host readiness} the ability of a pretrained world model to use an
external memory representation. After this definition, we use the more direct
term \emph{memory interface compatibility}.

\paragraph{Contributions.}
\begin{itemize}\setlength\itemsep{1pt}
\item We define WRITE, KEEP, and RECALL with one causal importance score: the
change in a fixed host's future-latent loss when an event group is removed.
\item We implement a label-free raw-image protocol with frozen DINOv2 patch
tokens, action-conditioned host surprise, delayed group verification, event
versioning, hysteresis, and bounded retrieval.
\item We report trajectory-disjoint results for 18 cells, 10 environments, and
four families, including reset, shuffled, random, recent-only, budget, delay,
rollout, deletion, and action-sequence controls.
\item We separate DINO-feature breadth from released DINO-WM validation and
show that storage quality and memory interface compatibility are distinct.
\end{itemize}

\section{Related Work}

\paragraph{Predictive world models.}
Joint-embedding predictive architectures learn by predicting latent targets
rather than pixels \citep{r_jepa}. DINO-WM predicts frozen DINO features and
uses them for planning \citep{r_dinowm}; LeWM trains an end-to-end latent world
model \citep{r_lewm}. Causal-JEPA studies object-level latent masking
\citep{r_causaljepa}. These methods learn predictive representations inside a
configured context. We study an external state that persists after an
observation leaves that context.

\paragraph{Learned memory.}
Recurrent networks, state-space models, and slot memories provide different
ways to compress history \citep{r_gru,r_lstm,r_mamba,r_slotssm,r_gsa}.
SlotMemory uses object-centered key--value storage \citep{r_slotmemory}.
Titans writes to neural memory using surprise \citep{r_titans}, and test-time
training adapts computation online \citep{r_ttt}. In \cem, surprise only
proposes an event. Retention is trained against the effect of removing the
event from a fixed external host.

\paragraph{Memory evaluation.}
RULER and LongMemEval show that nominal context length can overstate usable
memory \citep{r_ruler,r_longmemeval}. ``On Memory'' argues that memory should
be evaluated through downstream use \citep{r_onmemory}. We use a
\emph{controlled memory evaluation}: full, reset, recent-only, shuffled, and
random memories share the same host, data, and prediction target. Probe
accuracy is secondary to host prediction loss and deletion tests.

\paragraph{Surprise and predictive processing.}
Predictive coding treats mismatch between expected and observed activity as a
learning signal \citep{r_predictivecoding}; active inference extends prediction
error to action and belief updating \citep{r_activeinference}. CEM uses a
narrower engineering role: surprise proposes a write, while delayed future
prediction and hard deletion decide whether the event remains useful.

\paragraph{Interventional representation learning.}
Interventional causal representation learning uses controlled changes to
identify latent factors under explicit assumptions \citep{r_interventionalcrl}.
CEM instead intervenes on the memory input to an already fixed predictor. Its
causal statement is therefore local: deletion changed this host's future
prediction, not the environment's causal structure.

\section{Problem Definition}
\label{sec:problem}

Let a fixed host encode frame $x_t$ as $z_t$ and predict future latents from a
legal context of $K$ observations and actions. A memory store
$\mathcal M_t$ contains event groups
\[
m_i=(k_i,v_i,t_i,\widehat{\CE}_i,\ell_i),
\]
where $k_i$ and $v_i$ are a semantic key and value, $t_i$ is the observed
timestamp, $\widehat{\CE}_i$ is a predicted causal importance score, and
$\ell_i$ records whether the event is provisional, active, superseded,
rejected, or evicted.

For future evaluation set $\mathcal H_t$, define the measured group effect
\begin{equation}
\CE(G)=
\mathbb E_{u\in\mathcal H_t}\!\left[
\mathcal L_{\mathrm{future}}\!\left(\mathcal M_t\setminus G\right)
-\mathcal L_{\mathrm{future}}(\mathcal M_t)\right].
\label{eq:ce}
\end{equation}
The loss is the host's action-conditioned latent prediction error. We delete
adjacent frames as a group because one visual transition can span several
frames. Equation~\ref{eq:ce} measures an intervention on the memory input to a
fixed computation. It is not an estimate of a causal effect in the external
environment.

\begin{figure*}[t]
\centering
\includegraphics[width=.98\textwidth]{fig_d_architecture.pdf}
\caption{\textbf{CEM on an unmodified OGBench trajectory.} Four raw Cube-single
frames are selected by a fixed pixel-change plus DINO-change score, restricted
to retrieved test events. CEM discovers the $t{=}3$ change, rejects a later
low-value event, verifies the stored group by deletion, and recalls it at
$t{=}17$. No event time or label is supplied. Exact scores and hashes are in
\protect\path{generated_results/architecture_example_receipt.json}.}
\label{fig:architecture}
\end{figure*}

\section{Causal-Effect Memory}
\label{sec:method}

\subsection{WRITE: surprise and semantic change}

The fixed host predicts $\hat z_t$ from the legal context and action. Its
one-step surprise is
\[
s_t=\|\hat z_t-z_t\|_2^2.
\]
Separately, the method measures temporal change in frozen DINOv2 patch-token
features. Train-set quantiles set both thresholds. Adjacent threshold
crossings form a provisional event group, and the highest combined crossing
defines its timestamp. This procedure has no manually chosen interval, color
target, or saliency function.

\subsection{KEEP: delayed group verification}

A small estimator predicts the future-loss reduction from retaining a
provisional group. Training targets are computed periodically by comparing
the host with a singleton event group against the same host with no event.
The target is normalized by no-memory loss. Smooth-L1 regression provides
scale calibration, and a pairwise ranking loss trains ordering within each
query. A group is promoted only after the verification delay and only when its
predicted score exceeds a validation-selected threshold.

\subsection{Event versioning}

Event versioning means that semantically similar events share a key but keep
separate timestamps and scores. Train-only clustering of DINO event vectors
assigns keys. A new candidate does not remove an active event when it is
proposed. The old event remains available until the new candidate passes
verification by more than a hysteresis margin. If verification fails, the
store falls back to the old version. Capacity eviction removes the lowest
scored active event.

\subsection{RECALL: bounded routing}

For query $q_t$, the router scores active events by
\begin{equation}
r_i=q_t^\top k_i+\phi(t-t_i)+\nu(q_t,\mathrm{context})
+\lambda\widehat{\CE}_i ,
\label{eq:router}
\end{equation}
where $\phi$ is a learned age function and $\nu$ predicts whether the current
query needs memory. The router retrieves at most $k$ events. Verification
uses the measured singleton group loss after the future evaluation horizon.
We report this bounded retrieve-then-verify audit separately. The primary
prediction uses the routed event before this outcome is known, so its test
target never selects its own memory.

\begin{algorithm}[t]
\small
\caption{Online CEM with delayed verification}
\label{alg:cem}
\KwIn{Frozen host $h$, bounded store $\mathcal M$, stream $(x_t,a_t)$}
$z_t\leftarrow h_{\rm enc}(x_t)$; compute one-step surprise and DINO change\;
\If{train-calibrated thresholds cross}{append adjacent frames as a provisional group}
\ForEach{group whose verification delay expires}{
  predict $\widehat{\CE}(G)$; promote, reject, or supersede by hysteresis\;
}
evict the lowest-value active group if $|\mathcal M|$ exceeds budget\;
$C_t\leftarrow\operatorname{topk}_{G\in\mathcal M} r(G,q_t)$\;
\Return frozen-host prediction conditioned on $C_t$\;
\textbf{Audit only after the target arrives:} measure hard-delete $\CE(G)$\;
\end{algorithm}

\begin{table}[t]
\centering\footnotesize
\caption{\textbf{Signal-to-claim contract.} Each stage supports a narrower
claim than the next; only hard deletion through the frozen host is treated as
causal evidence.}
\label{tab:signal-contract}
\begin{tabularx}{\linewidth}{lYY}
\toprule
signal & controller role & licensed interpretation\\
\midrule
surprise + DINO change & DISCOVER & candidate, not importance\\
$\widehat{\CE}$ & STORE / replace & amortized predicted value\\
query/content/age/need & RECALL & bounded relevance ranking\\
hard group deletion & VERIFY / audit & effect on fixed-host prediction\\
\bottomrule
\end{tabularx}
\end{table}

\paragraph{Training/evaluation separation.}
Thresholds, key clustering, the effect estimator, and router are fit on train
trajectories and selected on validation trajectories. Test futures never
select the primary memory. The retrieve-then-verify result is a separate
post-hoc audit because it observes the realized future target.

\paragraph{Bounded cost.}
Surprise is a by-product of the host forward pass. Online KEEP and eviction use
the amortized score; true group deletion is sampled during calibration and
audit. The store, top-$k$ read, and version budget are fixed, so persistent
state does not enlarge the host's legal attention window.

\section{Experimental Protocol}
\label{sec:protocol}

\subsection{Raw OGBench breadth}

We use original 64$\times$64 render caches with 768 trajectories and 22 frames
per environment. The environments are PointMaze-large, PointMaze-teleport,
PointMaze-giant, Cube-single, Cube-double, Cube-triple, Puzzle-3x3, Scene,
AntMaze-large, and HumanoidMaze-large. PointMaze-large, PointMaze-teleport,
Cube-single, and Cube-triple use three optimization seeds; the other
environments use one breadth seed.

A fixed trajectory permutation creates 538 train, 115 validation, and 115 test
trajectories. A frozen DINOv2 ViT-S/14 produces a spatial pyramid of patch
tokens. A 96-dimensional PCA projection is fit on train trajectories only.
For each environment and seed, a finite-window action-conditioned predictor is
trained on train trajectories, selected by validation loss, and frozen before
CEM training. All 18 hosts beat latent persistence on test.

The model reads only `frames`, `actions`, and array timestamps. Other cache
keys are ignored. The validation split selects write thresholds, verification
delay, hysteresis, budget, and retrieval count. The test split reports
future-latent MSE by horizon, all memory controls, causal deletion, score
calibration, write and retrieval rates, budget and delay curves, and post-hoc
action-sequence identification.

\subsection{Separate controlled protocols}

The released DINO-WM Wall protocol freezes the published predictor,
action/proprioception encoders, and DINOv2 features. It evaluates event
versioning by relocating cached encoded events in time. This is a controlled
representation-level protocol, not raw OGBench breadth or planning.

The LeWM PushT protocol freezes the released encoder and predictor and tests a
six-way same-base visual binding. It measures the path from readable memory to
the predictor output. Semantic class labels are used only by post-hoc
readouts. These synthetic tasks diagnose overwrite, binding, and interface
failures; they are not the main breadth evidence.

\subsection{Decision anatomy}

\begin{figure*}[t]
\centering
\includegraphics[width=.92\textwidth]{cem_raw_event_timeline.pdf}
\caption{\textbf{One automatic decision trace.} On the same raw Cube rollout
used in Fig.~\ref{fig:architecture}, the $t{=}3$ event is verified and later
retrieved at $t{=}17$. A larger proposal at $t{=}9$ is rejected by delayed
verification. Proposal strength therefore does not determine retention.}
\label{fig:raw-timeline}
\end{figure*}

The trace separates DISCOVER from KEEP. The later event is more surprising,
but its predicted future value is below the active version's hysteresis
threshold. CEM retains the older event and retrieves it without a supplied
event boundary. The deterministic selection score and raw frame hashes are in
\path{generated_results/architecture_example_receipt.json}.

\section{Results}
\label{sec:results}

\subsection{Raw OGBench breadth}

\begin{figure*}[t]
\centering
\includegraphics[width=.94\textwidth]{fig_raw_breadth.pdf}
\caption{\textbf{Does CEM improve future prediction?} Per-environment
four-step prediction improvement, normalized to each control. CEM improves on
no memory in all ten environments, but recent-only is better in 15/18
optimization cells. Error bars show optimization-seed variation where
available; separators denote manipulation, navigation, puzzle, and scene.}
\label{fig:raw-prediction}
\end{figure*}

CEM lowers prediction error relative to no memory in every environment. The
mean reduction is $0.719\%$ and grows from $0.227\%$ at rollout step one to
$1.086\%$ at step four. Reset, shuffled, and matched-random memories do not
explain the gain. The stronger recent-only control does: CEM is slightly worse
on average and wins only 3/18 cells. The primary result is therefore useful
persistent state, not superior event selection. Raw MSE values and horizon
curves are reported in Appendix~\ref{app:raw-controls}.

\begin{figure*}[t]
\centering
\begin{minipage}[t]{.55\textwidth}
\includegraphics[width=\linewidth]{cem_raw_rollout_horizon.pdf}
\end{minipage}\hfill
\begin{minipage}[t]{.41\textwidth}
\includegraphics[width=\linewidth]{cem_raw_budget_pareto.pdf}
\end{minipage}
\caption{\textbf{When does memory help?} Left: normalized improvement grows
over the four-step rollout. Right: matched controls confirm that an arbitrary
residual does not help, while recent-only remains the strongest alternative.
Intervals treat the 18 optimization cells as observations.}
\label{fig:raw-horizon-controls}
\end{figure*}

\begin{figure*}[t]
\centering
\includegraphics[width=.94\textwidth]{fig_raw_causal.pdf}
\caption{\textbf{Does the learned score identify causal events?} Deletion
effects are normalized by each host's no-memory prediction error. On raw
OGBench, high-CE and matched-random groups have similar effects and vary by
environment. The released DINO-WM protocol yields a large selected-versus-random
contrast, but is a separate encoded-event study.}
\label{fig:raw-causal}
\end{figure*}

\subsection{Causal ranking is not yet reliable}

High-ranked deletion exceeds matched random deletion in 12/18 cells, but the
aggregate gap includes zero and held-out rank correlation is near zero. CEM
therefore discovers useful state more reliably than it orders individual
events. Calibration and action-sequence diagnostics are reported in
Appendix~\ref{app:raw-use}; neither supports a planning claim.

\subsection{Delayed verification controls the event store}

\begin{figure*}[t]
\centering
\includegraphics[width=.88\textwidth]{fig_discovery_progression.pdf}
\caption{\textbf{Which controller decisions matter?} Immediate surprise,
delayed causal-effect verification, the complete proxy controller, and
released DINO-WM are compared on false writes, correct overwrite, and retrieval
precision. Proxy bars average Cube and PointMaze; the official bar is a
separate Wall protocol.}
\label{fig:discovery-store}
\end{figure*}

Immediate surprise writes distractors. Delayed group verification sharply
reduces false writes and improves overwrite and retrieval. Event versions then
protect an active event while a replacement is tested. These task-level rates
motivate the full controller without relying on raw prediction-loss units.

\subsection{Released DINO-WM gives stronger controlled deletion evidence}

\begin{figure}[t]
\centering
\includegraphics[width=\linewidth]{fig_official_factorial.pdf}
\caption{\textbf{Released frozen DINO-WM Wall factorial} (three seeds).
The panels report task exposure, correct overwrite, normalized prediction-error
reduction, and selected-event deletion relative to the reset baseline. This
encoded controlled protocol is separate from raw OGBench breadth.}
\label{fig:official}
\end{figure}

On the released DINO-WM host, CEM exposes the event above chance, improves
prediction by $16.6\%$, and produces a $21.2\%$ normalized selected-event
deletion effect while random deletion is zero. This is clear host-use evidence,
but the protocol relocates cached encoded events and does not establish raw
breadth or planning.

\subsection{LeWM reveals memory interface incompatibility}

\begin{figure}[t]
\centering
\includegraphics[width=\linewidth]{fig_lewm_readiness.pdf}
\caption{\textbf{Frozen LeWM interface study} (six-way PushT binding, age 15,
three seeds). Left: task information available at the host output. Right:
relative prediction degradation for interfaces evaluated with the same strict
target. Non-comparable legacy objectives appear only in the left diagnostic.}
\label{fig:lewm}
\end{figure}

The memory state retains the six-way binding, but dense residual injection
loses it before the frozen predictor. Memory tokens and AdaLN expose the
binding, but severely degrade prediction. Strict hybrid and LoRA preserve
prediction but remain near chance. No tested point combines high exposure with
less than $10\%$ degradation. Storage quality and host-interface compatibility
are therefore separate.

\section{Interpretation}

\paragraph{Prediction and causal ordering are distinct.}
An imperfect event set can lower average prediction error even when its
within-query ranking is weak. CEM is consistent on the former endpoint and
mixed on the latter, which is why the raw breadth and deletion figures answer
separate questions.

\paragraph{Recency is the strongest alternative.}
Reset, shuffled, and random controls rule out a generic residual benefit.
Recent-only retains valid short-range trajectory state and remains stronger.
A decisive next benchmark must break local continuity with longer,
trajectory-disjoint gaps while preserving the same memory and retrieval
budget.

\paragraph{Breadth is directional, not equally precise.}
Four representative environments have three optimization seeds; six have one.
The ten-environment result establishes a shared direction over no memory, not
equal confidence across families. Full per-environment losses, uncertainty,
telemetry, and audit receipts remain in the appendix.

\section{Claim Hierarchy}

\begin{table}[t]
\centering\footnotesize
\caption{\textbf{Evidence is reported at four distinct levels.} A stronger
claim is not inferred from a weaker endpoint.}
\label{tab:claim-hierarchy}
\begin{tabularx}{\linewidth}{YcY}
\toprule
claim & status & evidence\\
\midrule
raw state improves on no memory & supported & 10/10 environments; matched controls\\
CEM beats local recency & unsupported & 3/18 cell wins\\
learned CE orders raw events & unresolved & deletion-gap interval includes zero\\
released host uses selected memory & supported & fail-closed audit + deletion\\
frozen LeWM is memory-compatible & unsupported & exposure--fidelity trade-off\\
\bottomrule
\end{tabularx}
\end{table}

\paragraph{Why the hierarchy matters.}
Prediction improvement establishes that retrieved history can be useful. It
does not establish that the written event is uniquely necessary, that the
learned score is calibrated, or that an arbitrary pretrained host can consume
the memory. The raw, official-host, and LeWM protocols isolate these questions
rather than averaging them into one headline metric.

\paragraph{What would change the conclusion.}
The raw selection claim would require CEM to outperform a capacity-matched
recent-only state across longer gaps and to show a positive high-versus-random
deletion interval. The host-compatibility claim would require a LeWM interface
above the exposure gate while remaining below the registered prediction
degradation bound. Neither condition is met here.

\paragraph{Why the official result remains separate.}
Released DINO-WM supplies the clearest selected-event deletion contrast, but it
uses relocated cached latent events. It validates the controller's causal-use
mechanism under a fixed released host; it cannot substitute for raw-event
breadth, automatic semantic discovery, or executed planning.

\section{Discussion and Limitations}

\paragraph{What the raw result establishes.}
Frozen DINO features, host surprise, and a small event store can improve
action-conditioned prediction without supplied event labels or times. The
increasing horizon gain and the shuffled/random controls support a genuine use
of past visual state. The result covers unmodified simulated environment
renderings and does not test real-world camera data.

\paragraph{What it does not establish.}
Recent-only state is stronger on average, CE-score calibration is low, and the
deletion-ordering interval includes zero. Trajectories contain only 22 frames
and four rollout steps are evaluated. Six environments have one breadth seed.
The PCA projection and memory adapter are environment-specific. These limits
rule out claims of unbounded recall, general causal discovery, or transfer.

\paragraph{Host dependence.}
Equation~\ref{eq:ce} can identify only information represented and used by the
fixed host. A blind host produces no useful score. Conversely, a memory can
retain information that the host interface removes, as in LeWM. Future work
should jointly train a memory-compatible host interface while retaining
controlled reset, shuffled, random, and recent-only comparisons.

\paragraph{Computation.}
Periodic group deletion requires extra host evaluations, and bounded
retrieve-then-verify scales with the number of routed candidates. The current
implementation limits this cost with amortized scoring and top-$k$ retrieval,
but does not report a standardized cross-host latency benchmark.

\section{Conclusion}

\cem\ extends a finite-window world model with label-free visual event memory
whose decisions are trained against future-latent prediction loss. On raw
OGBench renderings, it consistently improves over no memory and corruption
controls, with larger gains at longer rollout horizons. It does not yet beat a
recent-only extension or reliably rank event effects. Released DINO-WM and
LeWM diagnostics show why both event selection and memory interface
compatibility must be evaluated. These results provide a reproducible baseline
for raw visual event memory and a clear set of remaining failures.
